\solchapter{newton}

\sol{newton}{1} $R = v_0^2\sin2\theta/g$ is largest when $\sin2\theta = 1$, i.e.\ $\theta = 45^\circ$. From a cliff the projectile has extra time in the air while it falls the height $h$, and that extra time is worth more spent moving horizontally. So the optimum tilts toward flatter throws: less than $45^\circ$, approaching $0^\circ$ for a very high cliff.

\sol{newton}{2} Along the ramp, gravity pulls with $mg\sin\theta$; friction can supply at most $\mu_sN = \mu_smg\cos\theta$. Sliding starts when $\tan\theta>\mu_s$. The mass cancels. To measure $\mu_s$, tilt the ramp slowly until the block just starts to slide, and take $\mu_s = \tan\theta_c$.

\sol{newton}{3} For each mass along its direction of motion: $m_1a = m_1g - T$ and $m_2a = T - m_2g$. Adding,
\[
  a = \frac{m_1-m_2}{m_1+m_2}\,g, \qquad T = \frac{2m_1m_2}{m_1+m_2}\,g .
\]
As $m_2\to0$, $a\to g$ and $T\to0$: free fall. For $m_1 = m_2$, $a = 0$ and $T = mg$: balance.

\sol{newton}{4} The drop has mass $\tfrac43\pi r^3\rho = 4.2\times10^{-6}\un{kg}$ and area $\pi r^2 = 3.1\times10^{-6}\un{m^2}$. With $C_D\approx0.5$ for a sphere, $v_t = \sqrt{2mg/C_D\rho_{\text{air}}A}\approx6.6\un{m/s}$ (measured: about $6.5\un{m/s}$). Without air it would hit the ground from a kilometre up at $140\un{m/s}$. Air drag caps it at a few metres per second, and a $4\un{mg}$ drop at that speed carries less than $10^{-4}\un{J}$.

\sol{newton}{5} $GMm/r^2 = m\omega^2r$ with $\omega = 2\pi/T$ gives $T^2 = (4\pi^2/GM)r^3$. With $GM_\oplus = 3.99\times10^{14}\un{m^3/s^2}$ and $T = 86\,164\un{s}$, $r = (GMT^2/4\pi^2)^{1/3} = 4.22\times10^7\un{m}\approx6.6\,R_\oplus$, an altitude of about $35\,800\un{km}$.

\sol{newton}{6} Energy conservation gives $Mgh = \tfrac12Mv^2(1+\beta)$ with $\beta = I/MR^2$, so $v^2 = 2gh/(1+\beta)$, and the acceleration $g\sin\theta/(1+\beta)$ is constant. Mass and radius cancel; only the shape factor $\beta$ matters. Sphere ($\tfrac25$) beats disc ($\tfrac12$) beats hoop ($1$). The frictionless block has $\beta = 0$, wastes nothing on spinning, and wins.

\sol{newton}{7} Take, for example, payload $1$, total structure $10$, fuel $89$. Single stage: $\Delta v = u\ln(100/11) = 2.21u$. Two stages, each with $44.5$ fuel and $5$ structure: first burn $u\ln(100/55.5) = 0.59u$, then drop the empty stage and burn $u\ln(50.5/6) = 2.13u$; total $2.72u$. Staging wins because the rocket stops paying to accelerate empty tanks; the logarithm in the rocket equation punishes every kilogram of dead weight.

\sol{newton}{8} In the rotating frame the falling ball (downward speed $gt$) feels an eastward Coriolis acceleration $2\Omega gt$ at the equator. Integrating twice, $x = \tfrac13\Omega gt^3$. With $t = \sqrt{2h/g} = 4.5\un{s}$ and $\Omega = 7.3\times10^{-5}\un{s^{-1}}$, $x\approx2.2\un{cm}$ east. (The naive inertial estimate, extra speed $\Omega h$ times $t$, gives $3.3\un{cm}$; the factor $\tfrac23$ arises because the ball's eastward speed advantage over the ground shrinks as it falls.)

\sol{newton}{9} $v_{\text{esc}} = \sqrt{2GM/R} = \sqrt{2(4.9\times10^{12})/1.74\times10^6} = 2.4\un{km/s}$, so the survival threshold is about $0.4\un{km/s}$. Nitrogen at the Moon's $\sim390\un{K}$ daytime temperature moves at $\sqrt{3\kB T/m}\approx0.6\un{km/s}$: lost. For the Earth, $v_{\text{esc}}/6\approx1.9\un{km/s}$. In the hot upper atmosphere ($\sim1000\un{K}$), nitrogen moves at $0.9\un{km/s}$ and is kept, while hydrogen atoms move at about $5\un{km/s}$ and escape.

\sol{newton}{10} Build from the top. Suppose the top $k$ books have their combined centre of mass exactly over the edge of book $k+1$. Measure positions inward from that edge: the top $k$ books contribute centre of mass $0$, book $k+1$ contributes $L/2$, so the top $k+1$ books have centre of mass $(L/2)/(k+1) = L/2(k+1)$ inside the edge of book $k+1$. That is how far book $k+1$ can overhang book $k+2$. The top book overhangs the second by $L/2$, the second the third by $L/4$, and so on, for a total of $\tfrac L2H_n$ with $H_n = \sum_{k\le n}1/k$. Since $H_n\sim\ln n$ diverges, the overhang is unbounded. The top book clears the table when $H_n>2$, which first happens at $n = 4$ ($H_4 = 2.08$).

\sol{newton}{11} In the rotating frame, at angle $\theta$ from the bottom, the bead has potential energy $U = -mgR\cos\theta - \tfrac12m\Omega^2R^2\sin^2\theta$ (gravity plus centrifugal). Then $U' = mR\sin\theta\,(g - \Omega^2R\cos\theta)$, so equilibria are at $\theta = 0$, $\theta = \pi$, and $\cos\theta = g/\Omega^2R$ if $\Omega^2>g/R$. At the bottom, $U''(0) = mR(g - \Omega^2R)$: stable for $\Omega^2<g/R$, unstable above. At the new points, $U'' = m\Omega^2R^2\sin^2\theta>0$: stable. The bead must choose one side, though the setup is symmetric --- a pitchfork bifurcation.

\sol{newton}{12} Set $\boldsymbol\omega = (\epsilon_1, \Omega, \epsilon_3)$ and keep first order: $I_1\dot\epsilon_1 = (I_2-I_3)\Omega\epsilon_3$ and $I_3\dot\epsilon_3 = (I_1-I_2)\Omega\epsilon_1$. Then $\ddot\epsilon_1 = \lambda^2\epsilon_1$ with
\[
  \lambda^2 = \frac{(I_2-I_3)(I_1-I_2)}{I_1I_3}\,\Omega^2 > 0,
\]
since both factors are negative. The wobble grows as $e^{\lambda t}$. Repeating the calculation for rotation about axis $3$ gives $\lambda^2\propto(I_2 - I_3)(I_3 - I_1)<0$, and about axis $1$ gives $\lambda^2\propto(I_1-I_2)(I_3-I_1)<0$: oscillations, so those rotations are stable.
